\section{The actual Hodge spectrum and its exterior marking}
\label{sec:actual-hodge}

We now prove the spectral and form assertions of
Theorem~\ref{thm:main}. All norms in this section are the
geometric, unnormalized $L^2$ norms. On each smooth fibre,
the real Hodge Laplacian has its usual $H^2$ operator
domain and $H^1$ quadratic-form domain. Put
\[
 L_{\min}(p)=\min_i|\log|\tau_i(p)||,\qquad
 \mu_j(p)=\Eig_j(E_p,M),\quad j=1,2.
\]

\subsection{Multiplicative transfer and the complete count}

We recall the variational metric comparison in a form
which avoids differentiating the metric. If two metrics
$g,h$ on a closed six-manifold satisfy
$(1-\delta)g\le h\le(1+\delta)g$, their squared
$L^2$ norms in every degree differ by factors
$1+O(\delta)$. The differential $\dd$ and its graph
domain are unchanged. On $\Dom\dd/\ker\dd$, the quotient
norms differ by the same factors, since each is an
infimum over the same closed subspace. Min--max for
\[
 \frac{\|\dd u\|^2}{\|[u]\|^2}
\]
therefore compares every ordered coexact eigenvalue
multiplicatively. Exact positive eigenvalues in degree
two are coexact positive eigenvalues in degree one.
Taking the ordered union gives the full positive
degree-two comparison. This is the variational
comparison recalled in \cite[Section~6]{wuebben-hodge};
it needs no estimate of the difference of
codifferentials from a $C^0$ metric bound.

Fix $R,A$. Theorem~\ref{thm:ricci-flat-comparison}
compares the actual metric with the model by an
error whose limiting bound tends to zero as
$R\downarrow0$, $A\to\infty$. The model theorem
gives exactly two small positive eigenvalues,
relatively asymptotic to
$\Eig_j(E_\varepsilon,M_R)$. Its complementary
one-form and two-form gaps may depend on $R,A$.
The actual normalized scales give
\begin{equation}\label{eq:log-length-conversion}
 \log\frac{R}{2A\varepsilon_i(p)}
       =\frac13|\log|\tau_i(p)||+O_{R,A}(1).
\end{equation}
Since the link norm is $2\pi/3$, this implies
relative equivalence of $E_\varepsilon$ and
$E_p$, with error $O_{R,A}(L_{\min}^{-1})$.
This is a comparison of each nonnegative summand,
so it is independent of the condition number.
The quotient masses $M_R$ tend relatively to
$M$ under their common identification with $V/K$.
Generalized min--max now gives
\begin{equation}\label{eq:actual-relative-rates}
 \lambda_{2,j}^+(p)=(1+o(1))\mu_j(p),\qquad j=1,2,
\end{equation}
first at fixed $R,A$ and then in the ordered
limits. For each error tolerance the restriction
on $p$ applies throughout the full smooth locus.

Fixing one sufficiently accurate pair $R,A$
transfers the complementary gaps. A gap uniform
in all auxiliary choices is unnecessary.
Topology gives $b_2(X_p)=2$ throughout the smooth
family. Thus the full low two-form space consists
of two harmonic forms and two positive modes,
with no additional cap or exterior modes.

\subsection{Complex type and primitivity}

The parallel holomorphic volume form on a
Ricci-flat Kähler threefold, together with the
metric, gives a parallel $\operatorname{SU}(3)$-
equivariant isomorphism
$\Lambda^{0,1}\to\Lambda^{2,0}$.
It intertwines the Hodge Laplacians: the
connection terms agree under the parallel map,
and the curvature terms are the same holonomy
representation action. A nonzero positive
$(2,0)$ eigencomponent below the one-form gap
would therefore give a forbidden $(0,1)$
eigenform. Reality excludes the conjugate
component as well.

The Kähler contraction $\Lambda$ commutes
with the Laplacian. If $\Lambda u$ were
nonzero for a small positive eigenform, its
differential would be a nonzero one-form
eigenvector at the same eigenvalue. Similarly
$\delta u$ must vanish. Hence every small
positive form is primitive, of real type
$(1,1)$, and
$u=\lambda^{-1}\delta\dd u$ is coexact.
This applies to the entire space $S_p$.

\subsection{A fixed observer and capture of the whole cluster}

Choose the compact exterior $C_*$ from
Section~\ref{sec:introduction}, with its
central metric. Let $\Pi_*$ be its $L^2$
projection onto distributionally closed
two-forms with unrestricted tangential trace.
The first-order estimate and harmonic
cohomology map in Section~\ref{sec:model-energy}
make
\begin{equation}\label{eq:actual-fixed-observer}
 \mathcal A u=\pr_W^G[\Pi_*(u|_{C_*})]
\end{equation}
a bounded operation on exterior $L^2$.
The exterior identifications are fixed.
It annihilates every globally closed form
on $X_p$, whose restriction class lies in
$K$. On a central harmonic representative
$\eta_D$, it gives $\pr_W^G D$.
The cohomology coefficients can be read
from periods of the closed projection;
we do not integrate an arbitrary unclosed
$L^2$ form over cycles.

Compare the model metric $g$ with the
actual metric $h$ under the surgery maps.
If their inner products differ relatively
by $\delta$, their projections onto the
orthogonal complement of the same closed
space $\ker\dd$ differ in norm by $O(\delta)$.
In particular, for $u\perp_h\ker\dd$
and $v=P_gu$,
\begin{equation}\label{eq:metric-coexact-projection}
 \|u-v\|_g\le C\delta\|u\|_g,\qquad
 \dd v=\dd u.
\end{equation}
Indeed for $z\in\ker\dd$, the inner product
$\langle u,z\rangle_g$ is just the difference
of the two metric inner products and is
bounded by $C\delta\|u\|_g\|z\|_g$.
Taking a supremum over unit $z$ proves
the estimate. The subtracted form is closed.

If $u\in S_p$ is a unit vector, its energy
is at most $C/L_{\min}$. The model energy
of $v$ is exactly $\|\dd v\|_g^2$.
The complementary model gap bounds the
component of $v$ outside the model small
coexact space by
$C_{R,A}L_{\min}^{-1/2}$ in norm.
Thus $u$ is within $C\delta+o_{R,A}(1)$
of that entire model space. Section
\ref{sec:model-compactness} identifies it
with $W_R$ through its trials and their
mass convergence. On these trials the
fixed observer tends to
\[
 J_R:W_R\longrightarrow W,\qquad
                   D\longmapsto\pr_W^G D.
\]
Their central limits are closed representatives
in class $D$, so the assertion holds in
operator norm on the two-dimensional input
space. The transported mass tends to $M$
as $R\downarrow0$. Consequently
\begin{equation}\label{eq:actual-cluster-mass}
 \limsup_{p\to0}
 \sup_{\substack{u\in S_p\\\|u\|_h=1}}
 |M(\mathcal Au,\mathcal Au)-1|
       \le C\delta(R,A)+e_M(R),\qquad e_M(R)\to0.
\end{equation}
The observer constants remain bounded for
small $R$ because $C_*$ lies outside all
modified central regions and the quotient
masses tend to a positive form.

Taking the ordered limits proves uniform
mass convergence and injectivity of
$\mathcal A$ on $S_p$. Since both dimensions
are two, it is an isomorphism. In the
notation of \eqref{eq:observed-pair-intro},
polarization gives
\begin{equation}\label{eq:actual-observed-mass}
 (1-e(p))M\le N_p\le(1+e(p))M,\qquad e(p)\to0.
\end{equation}
No internal gap or eigenvalue-ratio assumption
has entered this argument.

\subsection{All surviving harmonic classes}

For $k\in K$, let $h_g(k)$ and $h_h(k)$
be model and actual harmonic representatives
of the corresponding smoothing class.
Their difference $e=h_h(k)-h_g(k)$ is exact.
Using their respective orthogonality to
exact forms gives
\[
 \begin{split}
 \|e\|_h^2
  &=-\langle h_g(k),e\rangle_h\\
  &=-\bigl(\langle h_g(k),e\rangle_h
                          -\langle h_g(k),e\rangle_g\bigr).
 \end{split}
\]
Hence $\|e\|_h\le C\delta\|h_g(k)\|_g$.
The model harmonic mass converges to
$G_R|_K$, which tends to $G|_K$.
The ordered limits prove the claimed
actual harmonic mass on all of $K$.
In particular both surviving classes,
not just the polarizing class, are retained.

\subsection{The weighted lower bound in the fixed marking}

The relative eigenvalues and observed mass
do not alone identify the marked energy.
We first prove a one-sided form inequality.
For an actual coexact $u$ take $v=P_gu$
as above. Since $u-v$ is globally closed,
\begin{equation}\label{eq:observer-projection-invariance}
 \mathcal Au=\mathcal Av.
\end{equation}
Let $C_R$ be the model exterior,
$\Pi_R$ its closed projection, and
$\mathcal A_R$ its quotient observer.
For small $R$, $C_*\subset C_R$.
Set $\zeta=\Pi_R(v|_{C_R})$. Its class
restricts unchanged to $C_*$, so
$\mathcal A\zeta=J_R\mathcal A_Rv$.
The first-order projection estimate gives
\begin{equation}\label{eq:observer-derivative-discrepancy}
 |\mathcal Av-J_R\mathcal A_Rv|
       \le C_R\|\dd v\|_g.
\end{equation}
Indeed $v-\zeta$ has $H^1(C_R)$ norm
bounded by that right side, and the fixed
observer is bounded on $L^2(C_*)$.
This derivative estimate is stronger than
mere convergence of the two observers.

The largest eigenvalue of $E_p$ in any
fixed Euclidean coordinates is
$O(L_{\min}^{-1})$. Also
$E_p(J_RD)=E_p(D)$, since the restriction
periods vanish on $K$. Combine
\eqref{eq:model-weighted-lower},
\eqref{eq:log-length-conversion} and
\eqref{eq:observer-derivative-discrepancy}
in the energy seminorm to obtain
\[
 \begin{split}
 E_p(\mathcal Au)^{1/2}
 &\le\bigl(1+o_{R,A}(1)+C_RL_{\min}^{-1/2}\bigr)
                                      \|\dd v\|_g\\
 &\le\bigl(1+C\delta(R,A)+o_{R,A}(1)\bigr)
                                      \|\dd u\|_h.
 \end{split}
\]
Every error is proportional to the input's
own energy norm. Taking ordered limits,
restricting to $S_p$ and transporting
by $\mathcal A$ gives
\begin{equation}\label{eq:actual-one-sided-energy}
 E_p\le(1+e_1(p))Q_p,\qquad e_1(p)\to0.
\end{equation}

\subsection{Relative control of the full energy form}

Under the observation isomorphism, the
generalized eigenvalues of $(Q_p,N_p)$
are the two actual eigenvalues. Therefore
\eqref{eq:actual-relative-rates} and
\eqref{eq:actual-observed-mass} imply
\begin{equation}\label{eq:actual-energy-determinant}
 \frac{\det Q_p}{\det E_p}
  =\frac{\lambda_{2,1}^+\lambda_{2,2}^+}{\mu_1\mu_2}
                       \frac{\det N_p}{\det M}
       \longrightarrow1.
\end{equation}
Each eigenvalue of
$B_p=E_p^{-1/2}Q_pE_p^{-1/2}$ is at least
$(1+e_1(p))^{-1}$ by
\eqref{eq:actual-one-sided-energy},
while their product tends to one by
\eqref{eq:actual-energy-determinant}.
The larger is at most
$\det(B_p)(1+e_1(p))$. Thus both tend to one,
which proves
\begin{equation}\label{eq:actual-observed-energy}
 (1-e(p))E_p\le Q_p\le(1+e(p))E_p.
\end{equation}
This determinant argument never bounds
$E_p^{-1}$ in a fixed norm.

For an actual orthonormal eigenbasis,
let $C_p$ have columns $\mathcal Au_j$.
Then
$C_p^{\mathsf T}N_pC_p=I$ and
$C_p^{\mathsf T}Q_pC_p=\Lambda_p$.
Conjugating \eqref{eq:actual-observed-mass}
and \eqref{eq:actual-observed-energy}
gives the two matrix estimates in
\eqref{eq:marked-matrix-intro}. In particular the
$(i,j)$ energy error is bounded by
$e(p)\sqrt{\lambda_{2,i}^+\lambda_{2,j}^+}$.
This remains meaningful for arbitrarily
separated eigenvalues.

Compactness gives a uniform error tending
to zero, without an explicit rate.
The eigenline and separated-path assertions
are proved next.
